% !TeX root = ../sections/02_exercises.tex
\subsection{1-C-2.}
\begin{exercise}
    数列 $\{a_n\}$ が $\alpha\in\mathbb{R}$ に収束し，
    数列 $\{b_n\}$ が $\beta\in\mathbb{R}$ に収束するとき，
    \[
        \lim_{n\to\infty}
        \frac{
            b_n a_1+b_{n-1}a_2+\cdots+a_{n-1}b_2+a_n b_1
        }{n}
        =
        \alpha\beta
    \]
    であることを証明せよ。
\end{exercise}
\begin{background}
    この問題では，収束列どうしを逆向きに掛け合わせた平均を扱う。

    分子は
    \[
        b_n a_1+b_{n-1}a_2+\cdots+a_{n-1}b_2+a_n b_1
    \]
    である。これをシグマ記号で書くと，
    \[
        \sum_{j=1}^{n}a_jb_{n+1-j}
    \]
    となる。

    したがって，示すべき極限は
    \[
        \lim_{n\to\infty}
        \frac{1}{n}
        \sum_{j=1}^{n}a_jb_{n+1-j}
        =
        \alpha\beta
    \]
    である。

    \begin{definition}[数列の収束]
		数列 $\{a_n\}$ が $\alpha\in\mathbb{R}$ に収束するとは，
        任意の $\epsilon>0$ に対して，ある自然数 $N$ が存在して，
        $n\geq N$ ならば
        \[
            |a_n-\alpha|<\epsilon
        \]
        が成り立つことをいう。
    \end{definition}

    \begin{remark}
        仮定より
        \[
            a_n\to\alpha,
            \qquad
            b_n\to\beta
        \]
        である。

        したがって，$a_j$ は最終的に $\alpha$ に近づき，
        $b_j$ は最終的に $\beta$ に近づく。

        ただし，この問題では積が
        \[
            a_jb_{n+1-j}
        \]
        という形になっているため，$a_j$ の添字 $j$ が大きいとき，
        $b_{n+1-j}$ の添字が必ず大きいとは限らない。
        そのため，各項を直接
        \[
            a_jb_{n+1-j}\approx \alpha\beta
        \]
        と見るよりも，和全体を分解して評価する方が扱いやすい。
    \end{remark}

    \begin{lemma}[収束列は有界]
		数列 $\{x_n\}$ が実数 $L$ に収束するならば，
        $\{x_n\}$ は有界である。
        すなわち，ある定数 $C>0$ が存在して，
        すべての $n\in\mathbb{N}$ に対して
        \[
            |x_n|\leq C
        \]
        が成り立つ。
    \end{lemma}

    \begin{lemma}[Cesaro の平均定理]
		数列 $\{x_n\}$ が $L\in\mathbb{R}$ に収束するならば，
        \[
            \frac{x_1+x_2+\cdots+x_n}{n}\to L
        \]
        である。
    \end{lemma}

    \begin{remark}
        有限和では，足す順番を入れ替えても値は変わらない。
        したがって
        \[
            \sum_{j=1}^{n}b_{n+1-j}
            =
            b_n+b_{n-1}+\cdots+b_1
            =
            \sum_{j=1}^{n}b_j
        \]
        である。

        よって，$b_n\to\beta$ であることから，Cesaro の平均定理を使えば
        \[
            \frac{1}{n}\sum_{j=1}^{n}b_{n+1-j}
            =
            \frac{1}{n}\sum_{j=1}^{n}b_j
            \to \beta
        \]
        が分かる。
    \end{remark}
\end{background}
\begin{strategy}
    まず，分子をシグマ記号で表す。
    \[
        b_n a_1+b_{n-1}a_2+\cdots+a_{n-1}b_2+a_n b_1
        =
        \sum_{j=1}^{n}a_jb_{n+1-j}
    \]
    である。

    したがって，示すべきことは
    \[
        \frac{1}{n}\sum_{j=1}^{n}a_jb_{n+1-j}
        \to
        \alpha\beta
    \]
    である。

    ここで，$a_j$ を極限値と誤差に分ける。
    すなわち
    \[
        a_j=\alpha+(a_j-\alpha)
    \]
    と書く。

    これを代入すると，
    \[
        \frac{1}{n}\sum_{j=1}^{n}a_jb_{n+1-j}
        =
        \frac{1}{n}\sum_{j=1}^{n}
        \{\alpha+(a_j-\alpha)\}b_{n+1-j}
    \]
    である。

    したがって，
    \[
        \frac{1}{n}\sum_{j=1}^{n}a_jb_{n+1-j}
        =
        \alpha
        \frac{1}{n}\sum_{j=1}^{n}b_{n+1-j}
        +
        \frac{1}{n}\sum_{j=1}^{n}(a_j-\alpha)b_{n+1-j}
    \]
    と分解できる。

    第一項について考える。
    有限和の順序を入れ替えると，
    \[
        \sum_{j=1}^{n}b_{n+1-j}
        =
        \sum_{j=1}^{n}b_j
    \]
    である。

    よって
    \[
        \alpha
        \frac{1}{n}\sum_{j=1}^{n}b_{n+1-j}
        =
        \alpha
        \frac{1}{n}\sum_{j=1}^{n}b_j
    \]
    となる。

    仮定より
    \[
        b_j\to\beta
    \]
    であるから，Cesaro の平均定理より
    \[
        \frac{1}{n}\sum_{j=1}^{n}b_j
        \to
        \beta
    \]
    である。

    したがって第一項は
    \[
        \alpha
        \frac{1}{n}\sum_{j=1}^{n}b_{n+1-j}
        \to
        \alpha\beta
    \]
    となる。

    次に，第二項
    \[
        \frac{1}{n}\sum_{j=1}^{n}(a_j-\alpha)b_{n+1-j}
    \]
    が $0$ に収束することを示す。

    仮定より
    \[
        b_n\to\beta
    \]
    であるから，収束列は有界である。
    したがって，ある定数 $C>0$ が存在して，
    すべての $m\in\mathbb{N}$ に対して
    \[
        |b_m|\leq C
    \]
    が成り立つ。

    このとき，すべての $j=1,2,\ldots,n$ に対して
    $n+1-j\in\mathbb{N}$ であるから，
    \[
        |b_{n+1-j}|\leq C
    \]
    が成り立つ。

    したがって，三角不等式より
    \[
        \left|
            \frac{1}{n}\sum_{j=1}^{n}(a_j-\alpha)b_{n+1-j}
        \right|
        \leq
        \frac{1}{n}\sum_{j=1}^{n}
        |a_j-\alpha|\,|b_{n+1-j}|
    \]
    である。

    さらに
    \[
        |b_{n+1-j}|\leq C
    \]
    より，
    \[
        \left|
            \frac{1}{n}\sum_{j=1}^{n}(a_j-\alpha)b_{n+1-j}
        \right|
        \leq
        C\frac{1}{n}\sum_{j=1}^{n}|a_j-\alpha|
    \]
    となる。

    ここで，仮定より
    \[
        a_j\to\alpha
    \]
    であるから，
    \[
        |a_j-\alpha|\to 0
    \]
    である。

    したがって，Cesaro の平均定理より
    \[
        \frac{1}{n}\sum_{j=1}^{n}|a_j-\alpha|
        \to
        0
    \]
    である。

    よって
    \[
        C\frac{1}{n}\sum_{j=1}^{n}|a_j-\alpha|
        \to
        0
    \]
    である。

    以上より，はさみうちの考え方から
    \[
        \frac{1}{n}\sum_{j=1}^{n}(a_j-\alpha)b_{n+1-j}
        \to
        0
    \]
    が従う。

    したがって，
    \[
        \frac{1}{n}\sum_{j=1}^{n}a_jb_{n+1-j}
        =
        \alpha
        \frac{1}{n}\sum_{j=1}^{n}b_{n+1-j}
        +
        \frac{1}{n}\sum_{j=1}^{n}(a_j-\alpha)b_{n+1-j}
    \]
    において，第一項は $\alpha\beta$ に収束し，
    第二項は $0$ に収束する。

    ゆえに
    \[
        \frac{1}{n}\sum_{j=1}^{n}a_jb_{n+1-j}
        \to
        \alpha\beta
    \]
    である。
\end{strategy}